% !TEX program = lualatex
% Tagged PDF/UA output: compile with LuaLaTeX.
\DocumentMetadata{pdfstandard=UA-2,pdfversion=2.0,lang=en-US,tagging=on}
\documentclass{./../../latex/handout}
\begin{document}
\thispagestyle{plain}
\newcommand{\mytitle}{Worksheet: Integration}
\myheader{\mytitle}

\vspace{-2em}
\section*{Rules of Integration}
\vspace{-1em}
Integration reverses differentiation: if $F'(x) = f(x)$, then $\int f(x)\, dx = F(x) + c$.
\begin{itemize}
  \item \underline{Power rule:} $\displaystyle \int x^n\, dx = \frac{x^{n+1}}{n+1} + c \quad (n \neq -1)$

  \textit{Example}. $\displaystyle \int x^4\, dx = \frac{x^5}{5} + c$
  \item \underline{Exponential rule:} $\displaystyle \int e^x\, dx = e^x + c$
  \item \underline{Log rule:} $\displaystyle \int \frac{1}{x}\, dx = \ln x + c \quad (x > 0)$
  \item \underline{Sum and multiple:} $\displaystyle \int [f(x) + g(x)]\, dx = \int f(x)\, dx + \int g(x)\, dx$, \quad $\displaystyle \int k f(x)\, dx = k \int f(x)\, dx$

  \textit{Example}. $\displaystyle \int (6x^2 + 2)\, dx = 2x^3 + 2x + c$
  \item \underline{Definite integral:} $\displaystyle \int_a^b f(x)\, dx = F(b) - F(a)$

  \textit{Example}. $\displaystyle \int_0^2 3x^2\, dx = \Big[x^3\Big]_0^2 = 8 - 0 = 8$
\end{itemize}

\textbf{Exercises}
\begin{enumerate}
  \item Find $\displaystyle \int (4x^3 - 2x + 1)\, dx$ \vspace{2.2cm}
  \item Find $\displaystyle \int 3e^x\, dx$
  \newpage
  \item Find $\displaystyle \int \frac{2}{x}\, dx$ \vspace{2.5cm}
  \item Evaluate $\displaystyle \int_0^2 (x^2 + 1)\, dx$ \vspace{2.5cm}
\end{enumerate}

\section*{Economic Applications}
\vspace{-0.5em}
Whenever we know a \textit{rate} (a marginal quantity, a flow over time) and want the \textit{total}, we integrate. An initial condition, such as the fixed cost or the starting capital stock, pins down the constant of integration.

\textbf{Exercise.} A firm's marginal cost is $MC(Q) = 4Q + 3$ and its fixed cost is 12.
\begin{enumerate}
  \item[(a)] Find the total cost function $C(Q)$. \vspace{3cm}
  \item[(b)] Using a definite integral, find the increase in total cost when output rises from $Q = 1$ to $Q = 3$.
\end{enumerate}

\newpage
\textbf{Exercise.} A firm invests at the rate $I(t) = 6t$ per year, and its capital stock at $t = 0$ is $K(0) = 50$.
\begin{enumerate}
  \item[(a)] Find the capital stock $K(t)$. \vspace{3cm}
  \item[(b)] How much capital is accumulated between $t = 1$ and $t = 2$? \vspace{4cm}
\end{enumerate}

\textbf{Exercise.} Demand for a good is $p = 12 - Q^2$ and the market price is $p = 3$. Find the quantity sold, then the consumer surplus (the area under the demand curve and above the price).

\end{document}
